% Source : cours/cm6.qmd, prêt de Léa. Le capital, pas les intérêts.
\begin{board}{6}{Crédit : création puis destruction de monnaie}
 {Le risque, les fonds propres et le coût du financement limitent le crédit. Exemple fictif.}
\stage{5.4}{1}{Un prêt : deux écritures}
\stage{16}{2}{Le paiement circule}
\stage{26.6}{3}{Le capital est remboursé}
\bank{5.4}{10}
\node[smalllabel,text width=3.5cm] at (3.2,8.9) {Actif};
\node[smalllabel,text width=3.5cm] at (7.5,8.9) {Passif};
\draw[drawline] (1.5,8.1)--(9.2,8.1) (5.35,6)--(5.35,8.1);
\node[smalllabel,text width=3.5cm] at (3.2,6.9) {Crédit\\$+\,1\,000$ €};
\node[smalllabel,text width=3.5cm] at (7.5,6.9) {Dépôt\\$+\,1\,000$ €};
\captext{5.4}{Le dépôt naît du prêt.}
\bank{13.5}{8.8}
\bank{18.4}{8.8}
\node[smalllabel,text width=3.5cm] at (13.5,7.7) {Banque A};
\node[smalllabel,text width=3.5cm] at (18.4,7.7) {Banque B};
\draw[flow] (13.5,12.4)--(18.4,12.4);
\noteat{16}{6.3}{Réserves entre banques}
\captext{16}{Le dépôt circule\\entre les comptes.}
\node[smalllabel,text width=3.5cm] at (24.5,11.7) {Crédit};
\node[smalllabel,text width=3.5cm] at (28.8,11.7) {Dépôt};
\draw[drawline] (22.8,10.8)--(30.5,10.8) (26.65,8.3)--(26.65,10.8);
\node[label,text width=3.5cm] at (24.5,9.5) {$-\,1\,000$ €};
\node[label,text width=3.5cm] at (28.8,9.5) {$-\,1\,000$ €};
\draw[amber,line width=2pt] (23.3,7)--(30,7);
\captext{26.6}{La créance et le dépôt\\diminuent ensemble.}
\end{board}
